\chapter*{Preface}
\addcontentsline{toc}{chapter}{Preface}

\section*{Motivation}

Mathematics in the twenty-first century is a vast, sprawling landscape. As the boundaries of the discipline have expanded, the specialization of its practitioners has intensified. Today, it is not uncommon for researchers in algebraic geometry, partial differential equations, and probability theory to speak languages so distinct that they struggle to read one another's work. Yet, the history of mathematics shows that its greatest leaps forward have occurred at the interfaces where these seemingly disparate disciplines collide.

The Wolf Prize in Mathematics, established in 1978 by Ricardo Wolf, is unique among international mathematical honors. Unlike the Fields Medal, which is restricted to mathematicians under the age of forty, the Wolf Prize recognizes lifetime achievement. Its laureates represent the architects of modern mathematics—thinkers who spent decades building the frameworks, proving the bedrock theorems, and establishing the deep, cross-disciplinary connections that define our current understanding.

This book, \textit{Wolf Prize in Mathematics: A Comprehensive Survey (1978--2024)}, was conceived to address a simple but pressing need: the lack of a single, unified resource that surveys the work of these 68 laureates. It is designed to act as a map of the modern mathematical landscape, outlining not only the historical biographies of the prize winners but also the precise mathematical structures they created. It is our hope that by reading about the roads built by Israel Gelfand, the modular curves analyzed by Andrew Wiles, the wavelets developed by Ingrid Daubechies, and the cryptographic algorithms designed by Adi Shamir, the reader will begin to see the underlying unity that binds all of mathematics together.

\section*{Scope and Structure}

This volume surveys all 68 Wolf Prize in Mathematics laureates from the inaugural awards in 1978 through the most recent awards in 2024. The book is organized chronologically, with each chapter dedicated to a single laureate (or shared laureates, treated individually). 

To ensure the book is accessible to readers at different stages of their mathematical development, each chapter is structured with clear pedagogical pathways. We have labeled sections with difficulty ratings:
\begin{itemize}
    \item \textbf{[B] Basic:} Accessible to advanced undergraduates with a standard background in linear algebra, calculus, and introductory abstract algebra.
    \item \textbf{[I] Intermediate:} Accessible to beginning graduate students or advanced undergraduates who have taken specialized courses (e.g., real analysis, algebraic topology).
    \item \textbf{[A] Advanced:} Aimed at graduate students and active researchers in the relevant fields.
\end{itemize}

Every chapter contains a biography and career overview, historical context explaining the origin of the core problems, detailed mathematical expositions with formal definitions and proof sketches, and a discussion of the laureate's wider impact on science and society. Frequently asked questions, glossaries, timelines, and graded exercises are provided to make this book a fully interactive learning experience.

\section*{Acknowledgments}

A project of this scale would have been impossible without the support and inspiration of many individuals and organizations. 

First and foremost, we express our profound gratitude to the Wolf Foundation in Israel for establishing and maintaining this prestigious award, and for their permission to reference the official citations and biographical details of the laureates.

We are deeply indebted to the global mathematical community. The theorems, definitions, and proofs detailed in these pages are the collective intellectual heritage of humanity, and we thank the countless mathematicians whose expository papers, lecture notes, and textbooks helped shape our presentation of these complex theories. 

We thank our colleagues and peers who reviewed draft chapters, corrected typographic and mathematical errors, and offered valuable insights on how to simplify dense proofs without losing mathematical rigor. Special thanks go to our graduate students, whose feedback on the exercises and glossary entries helped refine the book's pedagogical structure.

Finally, we thank our families and loved ones for their infinite patience and encouragement during the long nights and weekends spent compiling, writing, and editing this manuscript. 

\vspace{1.5cm}

\noindent \textbf{The Editors \& Compilers} \\
\textit{July 2026}
